% Provenance (never printed): the growth strand of ../economics_series.md (grade 8: growth rate, compound growth with
% powers, doubling time and the rule of 70, small differences over decades). Brief: sources/book1/BRIEFS.md, g8-04.
% Chart data: figdata/grade-8/04-growth-and-compounding (WDI NY.GDP.PCAP.KD 1960-2024; Maddison Project US; Brenland
% accounts; Kestria invented in the script).
\chapter{Growth and Compounding}\label{ch:g8:growth-and-compounding}

\begin{openingchart}{Real GDP per person in Argentina, the United States and South Korea,
1960--2024, in US dollars of 2015, on a times-ten scale. The three were chosen for their average
growth: about 1, 2 and 6 percent a year.}{World Bank, World Development Indicators,
NY.GDP.PCAP.KD}
\begin{tikzpicture}
\begin{semilogyaxis}[omtimesten, xmin=1960, xmax=2024, ymin=500, ymax=100000, ytickten={3,4,5},
  xtick={1960,1970,1980,1990,2000,2010,2020}, xlabel={Year},
  ylabel={Dollars of 2015 per person}, legend pos=south east]
\addplot table[x=year, y=slow] {figdata/grade-8/04-growth-and-compounding/out/three_speeds.csv};
\addplot table[x=year, y=mid] {figdata/grade-8/04-growth-and-compounding/out/three_speeds.csv};
\addplot table[x=year, y=fast] {figdata/grade-8/04-growth-and-compounding/out/three_speeds.csv};
\legend{Argentina ($0.86\,\%$ a year), United States ($1.99\,\%$), South Korea ($5.75\,\%$)}
\end{semilogyaxis}
\end{tikzpicture}
\end{openingchart}

% figdata: speeds (ARG 7397 -> 12774, 0.86; USA 18992 -> 66857, 1.99; KOR 1038 -> 37048, 5.75);
% facts (slow_ratio 1.7, mid_ratio 3.5, fast_ratio 35.7)
In 1960, a person in Argentina had an \omterm{def:g7:households:income}{income} of $7\,397$ dollars (in dollars of 2015), a person in
South Korea $1\,038$: Argentina was about seven times richer. In 2024, a person in Argentina had
$12\,774$ dollars, a person in South Korea $37\,048$. In 64 years, Argentina's \omterm{def:g6:living-standards:income-per-person}{income per person}
was multiplied by $1.7$, that of the United States by $3.5$, and South Korea's by $35.7$. The
yearly rates behind these numbers look close: $0.86\,\%$, $1.99\,\%$ and $5.75\,\%$. Over decades,
small differences in a yearly rate become enormous differences in \omterm{def:g7:households:income}{income}.

\begin{recall}
The \omterm{def:g6:reading-economic-charts:rate-of-change}{rate of change} from a starting value to a new value is $(\text{new} - \text{start}) \div
\text{start}$ (\cref{ch:g6:reading-economic-charts}). A power $a^n$ multiplies $a$ by itself $n$
times. Some poorer countries grow faster than the richest and catch up
(\cref{ch:g6:living-standards}).
\end{recall}

\section{The growth rate}

\begin{definition}[Growth rate]\label{def:g8:growth-and-compounding:growth-rate}
The \emph{growth rate}\index{growth rate} of a quantity over a year is its \omterm{def:g6:reading-economic-charts:rate-of-change}{rate of change} over
that year, in percent:
\[
\text{growth rate} = \frac{\text{value this year} - \text{value last year}}{\text{value last
year}} \times 100.
\]
\end{definition}

\begin{measure}[Real GDP per person]
Economic growth is measured with real \omterm{def:g8:gdp:gdp-per-person}{GDP per person}: GDP in the \omterm{def:g7:money:price}{prices} of a \omterm{def:g8:inflation-and-price-indices:base-year}{base year}
(\cref{ch:g8:inflation-and-price-indices}), divided by the population. Using \omterm{def:g8:inflation-and-price-indices:real-gdp}{real GDP} removes the
rise in \omterm{def:g7:money:price}{prices}; dividing by the population removes the growth that comes only from having more
people. Comparing countries also needs the same currency: the opening chart converts every
country's \omterm{def:g8:inflation-and-price-indices:real-gdp}{real GDP} into US dollars of 2015.
\end{measure}

\section{Growth on growth}

\begin{definition}[Compound growth]\label{def:g8:growth-and-compounding:compound-growth}
A quantity has \emph{compound growth}\index{compound growth} when each year it grows by a fixed
percentage of its value of the year before. After $n$ years at a rate of $g\,\%$, a starting value
$V$ becomes
\[
V \times \left(1 + \frac{g}{100}\right)^{n}.
\]
\end{definition}

Growth by a fixed percentage is growth on growth: the second year's growth applies to the first
year's growth too. A town of $100$ families that grows by $10\,\%$ a year has 110 families after one
year, $110 \times 1.1 = 121$ after two, $100 \times 1.1^{10} \approx 259.4$ after ten. Growing by a
fixed \emph{amount} of 10 families a year, it would have only 200.

\begin{omchart}{The same start, 100, grown for ten years by a fixed amount of 10 a year (simple
growth) and by $10\,\%$ a year (compound growth).}{Computed}
\begin{tikzpicture}
\begin{axis}[omrecord, ybar=0.5pt, bar width=5pt, xmin=-0.6, xmax=10.6, ymin=0, ymax=280,
  xtick={0,...,10}, xlabel={Year}, ylabel={Value}, height=0.45\linewidth, area legend,
  legend style={at={(0.04,0.96)}, anchor=north west}]
\addplot[fill=black!35, draw=none] table[x=year, y=simple] {figdata/grade-8/04-growth-and-compounding/out/compound.csv};
\addplot[fill=omDef, draw=none] table[x=year, y=compound] {figdata/grade-8/04-growth-and-compounding/out/compound.csv};
\legend{$+10$ a year, $+10\,\%$ a year}
\end{axis}
\end{tikzpicture}
\end{omchart}

\begin{definition}[Average annual growth rate]\label{def:g8:growth-and-compounding:average-growth-rate}
The \emph{average annual growth rate}\index{average annual growth rate} of a quantity over $n$
years is the fixed yearly rate that would take its starting value to its final value in $n$
years.
\end{definition}

The average rate of South Korea in the opening chart is $5.75\,\%$ because $1.0575^{64} \approx
35.7$: growing by $5.75\,\%$ each year for 64 years multiplies a value by about $35.7$, as South
Korea's \omterm{def:g6:living-standards:income-per-person}{income per person} was multiplied from 1960 to 2024. In no single year did it grow by
exactly $5.75\,\%$.

\begin{omfigure}
\includegraphics[width=0.34\linewidth]{images/book1/ai/g8-growth-door-frame.jpg}

{\small A child's height marked on a door frame, one line a year. Height grows by an amount each
year, not by a fixed percentage, and stops; \omterm{def:g6:living-standards:income-per-person}{income per person} can keep growing by a percentage.}
\end{omfigure}

\section{Doubling time}

\begin{definition}[Doubling time]\label{def:g8:growth-and-compounding:doubling-time}
The \emph{doubling time}\index{doubling time} of a quantity that grows at a fixed rate is the
number of years it takes to double.
\end{definition}

\begin{proposition}[Rule of 70]\label{prop:g8:growth-and-compounding:rule-of-70}
A quantity that grows by $g\,\%$ a year, with $g$ between about $0.5$ and $10$, doubles in about
$70 \div g$ years.
\end{proposition}
\admitted

The rule is proved in the volume for grades 10 to 12, with the logarithm. Here it can be checked
with powers: at $2\,\%$, $70 \div 2 = 35$ years, and indeed $1.02^{35} \approx 2.00$. At $7\,\%$,
it gives 10 years: $1.07^{10} \approx 1.97$.

\begin{model}{Steady growth}
\assumes \omterm{def:g6:living-standards:income-per-person}{Income per person} grows at the same yearly rate, year after year.
\predicts \omterm{def:g6:living-standards:income-per-person}{Income per person} doubles every $70 \div g$ years; on a \omterm{def:g6:reading-economic-charts:times-ten-scale}{times-ten scale}, its path is a
straight line.
\tested United States, 1870--2022: an average of $1.66\,\%$ a year, a doubling about every 42
years; the record stays close to the straight line of the chart below for a century and a half.
\breaks Countries catching up grow fast for some decades and then slow down (South Korea);
wars, crises and depressions (the United States after 1929); countries whose growth stalls for
decades (Argentina since 1960).
\end{model}

\begin{omchart}{GDP per person in the United States, 1820--2022, in dollars of 2011, on a
times-ten scale, with the steady path at the 1870--2022 average rate of $1.66\,\%$ a year
(dashed).}{Maddison Project Database}
\begin{tikzpicture}
\begin{semilogyaxis}[omtimesten, xmin=1820, xmax=2022, ymin=1000, ymax=100000, ytickten={3,4,5},
  xtick={1820,1870,1920,1970,2020}, xlabel={Year}, ylabel={Dollars of 2011 per person},
  legend pos=north west]
\addplot[omDef, thick] table[x=year, y=gdppc] {figdata/grade-8/04-growth-and-compounding/out/us_two_centuries.csv};
\addplot[black!70, dashed, thick, unbounded coords=jump] table[x=year, y=steady] {figdata/grade-8/04-growth-and-compounding/out/us_two_centuries.csv};
\legend{GDP per person, steady $1.66\,\%$ path}
\end{semilogyaxis}
\end{tikzpicture}
\end{omchart}

% figdata: facts (us_1870 4803, us_2022 58488, us_growth_1870 1.66, us_1929 11954, us_1933 8048)
The path is close to the dashed line, with one large dip: \omterm{def:g8:gdp:gdp-per-person}{GDP per person} fell from $11\,954$
dollars in 1929 to $8\,048$ in 1933, and came back close to the line only in the 1940s.

\section{Small differences, large gaps}

Two countries start with the same \omterm{def:g6:living-standards:income-per-person}{income per person}. One grows by $1\,\%$ a year, the other by
$3\,\%$. After 35 years, the first has multiplied its \omterm{def:g7:households:income}{income} by $1.01^{35} \approx 1.42$, the
second by $1.03^{35} \approx 2.81$: the second is twice as rich. After 70 years, the factors are
about $2.01$ and $7.92$. Over decades, a difference of one or two points in the yearly \omterm{def:g8:growth-and-compounding:growth-rate}{growth rate}
adds up to a larger difference in \omterm{def:g7:households:income}{income} than most events of a single year.

\section{Exercises}

\begin{exercise}[$\star$]\label{exo:g8:growth-and-compounding:1}
Use the text. What was \omterm{def:g6:living-standards:income-per-person}{income per person} in South Korea and in Argentina in 1960 and in 2024?
\end{exercise}

\begin{exercise}[$\star$]\label{exo:g8:growth-and-compounding:2}
Real \omterm{def:g8:gdp:gdp-per-person}{GDP per person} rises from $20\,000$ to $20\,600$ dalens in a year. Compute the \omterm{def:g8:growth-and-compounding:growth-rate}{growth rate}.
\end{exercise}

\begin{exercise}[$\star$]\label{exo:g8:growth-and-compounding:3}
A value of $1\,000$ grows by $5\,\%$ a year for 3 years. Compute its value after each year.
\end{exercise}

\begin{exercise}[$\star$]\label{exo:g8:growth-and-compounding:4}
Using the rule of 70, how long does a quantity growing by $7\,\%$ a year take to double?
\end{exercise}

\begin{exercise}[$\star\star$]\label{exo:g8:growth-and-compounding:5}
Compute $1.02^{35}$ with a calculator. What does the result say about growth at $2\,\%$ a year?
\end{exercise}

\begin{exercise}[$\star\star$]\label{exo:g8:growth-and-compounding:6}
\omterm{def:g6:living-standards:income-per-person}{Income per person} is $10\,000$ dalens. Compute it after 20 years of growth at $1\,\%$ a year, and at
$3\,\%$ a year.
\end{exercise}

\begin{exercise}[$\star\star$]\label{exo:g8:growth-and-compounding:7}
South Korea's \omterm{def:g6:living-standards:income-per-person}{income per person} was multiplied by $35.7$ in 64 years. Check with a calculator that
$1.0575^{64}$ is close to $35.7$.
\end{exercise}

\begin{exercise}[$\star\star$]\label{exo:g8:growth-and-compounding:8}
A country's \omterm{def:g6:living-standards:income-per-person}{income per person} is $2\,000$ dollars and grows by $5\,\%$ a year. (a) Use the rule of 70
to estimate it after 28 years. (b) Compute it with a power.
\end{exercise}

\begin{exercise}[$\star\star$]\label{exo:g8:growth-and-compounding:9}
Use the text on the United States in 1929--1933. By what percentage did \omterm{def:g8:gdp:gdp-per-person}{GDP per person} fall? Does this
fit the model ``Steady growth''?
\end{exercise}

\begin{exercise}[$\star\star\star$]\label{exo:g8:growth-and-compounding:10}
A poorer country has an \omterm{def:g6:living-standards:income-per-person}{income per person} of $5\,000$ dollars growing by $6\,\%$ a year; a richer one
has $20\,000$ dollars growing by $2\,\%$ a year. Make a table every 10 years, then year by year near
the crossing, to find after how many years the poorer country passes the richer one.
\end{exercise}

\begin{exercise}[$\star\star\star$]\label{exo:g8:growth-and-compounding:11}
A country's \omterm{def:g8:growth-and-compounding:growth-rate}{growth rate} falls from $4\,\%$ to $2\,\%$ a year. How does its \omterm{def:g8:growth-and-compounding:doubling-time}{doubling time} change?
\end{exercise}

\begin{exercise}[$\star\star\star$]\label{exo:g8:growth-and-compounding:12}
The rule of 70 tells how fast a country doubles its \omterm{def:g7:households:income}{income} at a given rate. Explain what it cannot
tell about why one country grows faster than another.
\end{exercise}

\section{Problem: Two Neighbours, Sixty Years}

\begin{problem}[{Weekend problem --- real GDP per person of Brenland and Kestria, their growth
rates, and when Kestria would catch up}]
\label{pb:g8:growth-and-compounding:1}
% figdata: neighbours (2000 24680 10400; 2005 26270 13040; 2010 26170 15900; 2015 28390 19300; 2020 28370 22730;
% 2025 31470 27700). Kestria 2025 = 30000 kests (g8-02 ex. 6, g9-05) x 1.242 / 1.345, from the frozen package.
Everything in this problem is invented. The table gives real \omterm{def:g8:gdp:gdp-per-person}{GDP per person} of Brenland and of its
neighbour Kestria, in dalens of 2015 (Kestria's converted into dalens at the \omterm{def:g7:markets-and-prices:market}{market} exchange rate).
\begin{center}
\begin{tabular}{lrrrrrr}
\toprule
Year & 2000 & 2005 & 2010 & 2015 & 2020 & 2025 \\
\midrule
Brenland & 24\,680 & 26\,270 & 26\,170 & 28\,390 & 28\,370 & 31\,470 \\
Kestria & 10\,400 & 13\,040 & 15\,900 & 19\,300 & 22\,730 & 27\,700 \\
\bottomrule
\end{tabular}
\end{center}

\medskip
\noindent\textbf{Part I --- The record.}
\begin{enumerate}
  \item How many times richer per person was Brenland than Kestria in 2000? In 2025?
  \item By what percentage did each country's real \omterm{def:g8:gdp:gdp-per-person}{GDP per person} grow from 2000 to 2025?
  \item In which five-year periods did Brenland's real \omterm{def:g8:gdp:gdp-per-person}{GDP per person} fall?
  \item Compute Kestria's growth from 2020 to 2025, in percent.
\end{enumerate}

\medskip
\noindent\textbf{Part II --- Growth rates.}
\begin{enumerate}[resume]
  \item Compute $1.01^{25}$ and $1.04^{25}$. Which is close to Brenland's growth factor from 2000 to
        2025, and which to Kestria's?
  \item Using the rule of 70, what is the \omterm{def:g8:growth-and-compounding:doubling-time}{doubling time} of each country at these rates?
  \item At these rates, in which year would Brenland's 2025 \omterm{def:g6:living-standards:income-per-person}{income per person} double? And Kestria's?
  \item Explain why a country far behind can grow faster for a long time.
\end{enumerate}

\medskip
\noindent\textbf{Part III --- Projections.} From 2025, assume Brenland grows by $1\,\%$ and Kestria by
$4\,\%$ a year.
\begin{enumerate}[resume]
  \item Compute both countries' real \omterm{def:g8:gdp:gdp-per-person}{GDP per person} in 2035, to the nearest 10 dalens. Which is
        higher?
  \item Compute both in 2030. Between which two of the years 2025, 2030 and 2035 does Kestria pass
        Brenland?
  \item Find, year by year, the first year in which Kestria's real \omterm{def:g8:gdp:gdp-per-person}{GDP per person} is above
        Brenland's.
  \item Name one reason why these projections may turn out wrong.
\end{enumerate}
\end{problem}
